The maintenance-extended model of Chapter~\ref{ch:maintenance} captures the
basic logic of an aircraft-hours check, but not the richer nested maintenance
structure considered in this thesis. We represent that structure by four check
levels, denoted $A$, $B$, $C$, and $D$, with distinct trigger mechanisms,
durations, and reset semantics. This chapter extends the compact assignment
formulation to this A/B/C/D regime while
retaining the key modelling principle of Chapter~\ref{ch:maintenance}: a
schedule remains compact and assignment-based, with maintenance decisions
represented through binary indicators rather than explicit route enumeration.

The extension is not merely a bookkeeping exercise. If the model treats all
checks as equivalent or ignores the fact that a heavier check covers lighter
ones, it can certify schedules that violate the required maintenance cadence or
overstate the amount of slack available between inspections. The hierarchy is
therefore introduced through a set of dominance relationships and check-reset
indicators that preserve the correct operational semantics while keeping the
model tractable.

\section{Application representation of the hierarchy}
\label{sec:hierarchy_application}

In the application, the hierarchy is represented directly in the JSON instance
and then loaded by the Python model builder. The instance contains the aircraft
set, flight legs and their departure and arrival times, initial aircraft
locations, initial maintenance exposure, maintenance thresholds, maintenance
durations, station capacities, and the flight--aircraft cost matrix. The keys
for the four check levels are interpreted consistently throughout the
application: A and B thresholds are measured in accumulated flight minutes,
whereas C and D thresholds are measured in elapsed calendar days. Maintenance
durations remain expressed in minutes because they consume time in the daily
aircraft timeline.

This data flow creates a clear separation between instance data and model
logic. The input file specifies the regulatory and operational state at the
start of the planning horizon; \texttt{MILP\_Sheduler} constructs the binary
variables and maintenance constraints; and the resulting solution reports
aircraft assignments and maintenance decisions for the selected horizon. The
same instance format is consumed by \texttt{Scheduler}, the heuristic engine,
which maintains the corresponding A/B flight-minute counters and C/D day
counters while constructing routes. Consequently, exact and heuristic runs
start from the same fleet data and apply the same dominance order. This common
input is necessary but not sufficient for strict objective comparison:
Chapter~\ref{ch:experiments} separately validates whether an MILP assignment
can be replayed under the timeline simulator's initial-maintenance and ferry
semantics.

At application level, a maintenance check is not entered as an arbitrary
independent event after optimization. Its type, location, duration, and effect
on the aircraft state determine which subsequent flights are feasible. A
heavier check is therefore propagated through the hierarchy both in the MILP
variables and in the simulator: a D-check resets all four obligations, a
C-check resets A, B, and C, a B-check resets A and B, and an A-check resets
only A. This shared interpretation is what allows the application to produce
not only an assignment matrix, but an operationally interpretable aircraft
schedule.

\section{Check types and regulatory hierarchy}
\label{sec:hierarchy_regime}

Let
\begin{equation}
    \mathcal{C} = \{A,B,C,D\}
\end{equation}
be the ordered set of maintenance classes, with dominance order
\begin{equation}
    D \succ C \succ B \succ A.
\end{equation}
A check of type $c$ is said to satisfy all lighter checks $c' \prec c$ in the
same period. This means that a D-check resets the elapsed-day requirement for
C, and simultaneously satisfies the $B$- and $A$-check conditions; similarly, a
C-check covers the lighter levels below it. The hierarchy therefore allows the
model to represent a single maintenance event as a check of level $c$ while
implicitly carrying the obligations of all lower levels that are dominated by
that event.

For each aircraft $j \in \Pc$, day $d \in \D$, and check type $c \in \C$, define
\begin{equation}
    y_{jdc} =
    \begin{cases}
        1, & \text{if aircraft } j \text{ undergoes a check of type exactly } c \text{ on day } d,\\
        0, & \text{otherwise.}
    \end{cases}
\end{equation}
The exact-type variable is complemented by the hierarchy indicator
\begin{equation}
    \gamma_{jdc} =
    \begin{cases}
        1, & \text{if aircraft } j \text{ undergoes a check of level at least } c \text{ on day } d,\\
        0, & \text{otherwise.}
    \end{cases}
\end{equation}
The indicator is a cumulative summary of the maintenance event: it is equal to
one whenever a check of level at least $c$ is scheduled for aircraft $j$ on day
$d$. In the compact model, we enforce the relation
\begin{equation}
    \gamma_{jdc} = \sum_{c' \succeq c} y_{jdc'},
    \qquad \forall j \in \Pc,\ d \in \D,\ c\in\C,
\end{equation}
which ensures that a heavier check automatically implies the lighter reset state
for that same day.

This construction is essential. It prevents a schedule from declaring an
aircraft as meeting the $A$-check requirement while it is actually undergoing a
heavier event that should satisfy the same obligation. The hierarchy is not a
separate planning layer; rather, it is the mechanism through which the model
maintains internally consistent reset semantics across all maintenance classes.

\section{Hierarchy constraints}
\label{sec:hierarchy_constraints}

The dominance structure is encoded through a small set of linear constraints.
First, at most one check level may be performed for a given aircraft on a given
calendar day:
\begin{equation}
    \sum_{c \in \C} y_{jdc} \leq 1,
    \qquad \forall j \in \Pc,\ d \in \D.
    \label{eq:hier_day_unique}
\end{equation}
This operational rule is consistent with the daily maintenance regime in which an
aircraft is either free, or engaged in a single maintenance event, but not both.

Second, the hierarchy indicator must respect subsumption. If $c' \succ c$, then
\begin{equation}
    \gamma_{jdc} \geq \gamma_{jdc'},
    \qquad \forall j \in \Pc,\ d \in \D,\ c \prec c'.
    \label{eq:hier_subsumes}
\end{equation}
The inequality states that a heavier event on day $d$ implies the aircraft is
also counted as having undergone all lighter checks on that day. Consequently,
 the same maintenance action simultaneously satisfies the reset condition for all
 dominated check types.

Finally, the exact-level and cumulative-level variables are tied together by
\begin{equation}
    \gamma_{jdA} = y_{jdA} + y_{jdB} + y_{jdC} + y_{jdD},
\end{equation}
\begin{equation}
    \gamma_{jdB} = y_{jdB} + y_{jdC} + y_{jdD},
\end{equation}
\begin{equation}
    \gamma_{jdC} = y_{jdC} + y_{jdD},
\end{equation}
\begin{equation}
    \gamma_{jdD} = y_{jdD},
\end{equation}
which are the instantiated forms of the general relation above. These equalities
are computationally convenient because they admit a sparse and transparent
definition of the maintenance state without introducing unnecessarily large sets
of independent variables.

\section{Flight-hour checks: types A and B}
\label{sec:hierarchy_ab}

Types $A$ and $B$ are triggered by cumulative block time. Let $T_{\max}^{c}$
denote the allowable accumulation limit for check type $c \in \{A,B\}$.
The corrected endpoint-anchored formulation of Chapter~\ref{ch:maintenance}
extends directly to the multi-type case by replacing the single reset variable
$y_{jd}$ with the relevant hierarchical indicator $\gamma_{jdc}$.

For a candidate interval $(d,d')$ and a check type $c \in \{A,B\}$,
\begin{align}
    \sum_{i \in F_{d,d'}} x_{ij}\,\tilde{t}_i
    &\leq T_{\max}^{c} + M_{dd'}\,\gamma_{jdc}
    + M_{dd'}\sum_{r=d+1}^{d'-1}\gamma_{jrc},
    \label{eq:c13a_hier} \\
    \sum_{i \in F_{d,d'}} x_{ij}\,\tilde{t}_i
    &\leq T_{\max}^{c} + M_{dd'}\,\gamma_{jd'c}
    + M_{dd'}\sum_{r=d+1}^{d'-1}\gamma_{jrc},
    \label{eq:c13b_hier}
\end{align}
where $F_{d,d'}$ is the set of flights departing between the two boundary days.
The logic of the correction is retained: a reset at either endpoint or any
intermediate day absorbs the flight-hour exposure, and the model does not
falsely permit the cumulative hours to exceed the threshold when a required check
should have intervened.

This generalization is the key structural step from a single maintenance class to
 a modelled four-level regime. A type-A or type-B reset is not only activated by a
 same-level event; it may also be triggered by a heavier check, because the
 heavier event satisfies the lighter requirement through $\gamma_{jdc}$. This is
 precisely what preserves the dominance ordering in the cumulative-hour
 constraints.

\subsection{Implemented spacing regime}

\textit{The implemented model does not impose an additional calendar-spacing
rule on the A/B flight-hour checks. Their compliance is governed by the
endpoint-anchored cumulative-hour rows above, including their pre-horizon
conditions. Calendar-day spacing is reserved for C/D checks, as specified in
the next section. This separation avoids introducing an unreported additional
restriction on the A/B feasible region.}

\section{Calendar-day checks: types C and D}
\label{sec:hierarchy_cd}

Types $C$ and $D$ are governed by elapsed calendar days rather than cumulative
 flight hours. For each such check type $c\in\{C,D\}$, let
 $d_{\max}^{c}$ denote the maximum admissible interval between two consecutive
 checks of level at least $c$. The spacing constraint is then
\begin{equation}
    \sum_{r=d}^{d+d_{\max}^{c}-1}\gamma_{jrc} \geq 1,
    \qquad \forall j \in \Pc,\ d \in \D,\ c\in\{C,D\}.
    \label{eq:spacing_cd}
\end{equation}
This rule says that in every window of length $d_{\max}^{c}$, aircraft $j$ must
experience at least one check of level at least $c$. Equivalently, there can be
no unbroken period longer than the maximum allowable interval without a qualifying
 maintenance intervention.

When a check type has a duration of multiple model days, a check that begins on
day $d$ occupies a contiguous maintenance window of length $\ell_c$, so the
aircraft is grounded on the days
\begin{equation}
    d,\ d+1,\ \ldots,\ d+\ell_c-1.
\end{equation}
To distinguish a new check from its continuation, set
$\gamma_{j,0,c}=0$. A newly active hierarchy indicator is propagated to the
following days by
\begin{equation}
    \gamma_{j,d+k,c} \geq \gamma_{jdc}-\gamma_{j,d-1,c},
    \qquad \forall j\in\Pc,\ d\in\D,\ c\in\C:\ell_c>1,\
    k=1,\ldots,\ell_c-1,
    \label{eq:multi_day_check}
\end{equation}
for indices that remain inside the planning horizon. The subtraction prevents
a continuation day from starting a fresh $\ell_c$-day chain. In the Pyomo
implementation, C14b imposes the equivalent binary start condition with a
big-$M$ row and uses the hierarchy variable \texttt{m.mega}.

To preserve feasibility, the model prevents a flight from being assigned to an
aircraft when its departure overlaps the active maintenance timing window. Let
$\mathcal W_{i'jdc}$ denote the candidate flights whose departure time and airport
fall within the duration window associated with maintenance triggered by flight
$i'$. \textit{The blocking condition is tied to the selected maintenance trigger,}
\begin{equation}
    z_{i'jdc}+x_{ij}\leq1,
    \qquad \forall j\in\Pc,\ c\in\C,\ d\in\D,\ i'\in\F,
    \ i\in\mathcal W_{i'jdc},
    \label{eq:no_flight_during_maintenance}
\end{equation}
\textit{or, in the baseline aggregate implementation, by the equivalent
C8 row that sums the blocked assignments and relaxes it only when
$z_{i'jdc}=0$. Thus a flight is blocked only when it overlaps the actual
maintenance interval created by its chosen trigger; no separate C15
calendar-day blocking row is used.}

\section{Pre-horizon state accounting}
\label{sec:hierarchy_prehorizon}

The maintenance logic remains valid only if the aircraft state at the beginning
of the planning horizon is treated consistently. For each aircraft $j$, let
\begin{itemize}
    \item $\Phi_j^{A}$ and $\Phi_j^{B}$ be the accumulated flight hours since the
    preceding A/B check,
    \item $\Psi_j^{C}$ and $\Psi_j^{D}$ be the elapsed calendar days since the
    preceding C/D check.
\end{itemize}
These values are not optional modelling extras; they are required to preserve the
continuity of regulatory compliance from the previous maintenance cycle into the
planning horizon.

Without this carry-forward information, the model implicitly assumes that every
aircraft begins the horizon in a fully reset state, which is unrealistic and may
 cause the generated schedule to violate the maintenance requirements within the
 first few days of operation. To avoid this, the initial sub-interval is adjusted
 so that the effective exposure is measured relative to the entry state, rather
 than assuming a zero-age aircraft at day one. The same principle applies to both
 hour-based and day-based checks: the first relevant interval in the horizon is
 shifted by the pre-horizon accumulated value, analogous to the correction already
 described in Chapter~\ref{ch:maintenance} for the single-check model.

This treatment is critical in computational practice because a schedule can appear
 feasible on the basis of a clean horizon start but still be invalid once the
 carry-over exposure is restored. The refined formulation therefore maintains
 operational continuity across planning periods, which is necessary when the
 observed schedule is intended to be used in a realistic fleet-operations setting.

For the A/B hierarchy, the opening flight-hour condition is explicitly
\begin{equation}
    \sum_{i\in F_{0,d'}}\tilde t_i x_{ij}
    \leq (T_{\max}^{c}-\Phi_j^c)
    +M_{0d'}\sum_{r=1}^{d'-1}\gamma_{jrc}
    +M_{0d'}\gamma_{jd'c},
    \qquad \forall j\in\Pc,\ c\in\{A,B\},\ d'\in\D.
    \label{eq:c13b_initial_hier}
\end{equation}
For C/D checks, the corresponding opening condition is a deadline on the first
qualifying $\gamma_{jdc}$ event, computed from the elapsed pre-horizon days.

\section{Model size and operational implications}
\label{sec:hierarchy_impact}

The full A/B/C/D extension increases the number of maintenance variables and
constraints relative to the basic maintenance model, but it does so in a controlled
and structured way. The exact-type variables $y_{jdc}$ are defined for each
 aircraft, day, and check level, so the additional binary dimension is of order
 $O(|\Pc|\,|\D|\,|\C|)$. The hierarchy indicator variables $\gamma_{jdc}$ add
 the same order of magnitude, while the spacing, dominance, and heavy-check
 constraints contribute a comparable set of linear rows. The extended formulation
 therefore remains compact and assignment-based; it does not enumerate complete
 aircraft routes or repair schedules explicitly.

\begin{proposition}
    The addition of the full $\{A,B,C,D\}$ hierarchy to the maintenance-extended
    compact model increases the number of maintenance-related binary variables by a
    factor bounded by $|\mathcal{C}|=4$ and introduces a linear constraint set of
    comparable order. The model remains polynomial in the size of the daily
    assignment network and does not require route enumeration.
\end{proposition}

The operational value of the extension is substantial. A compact formulation that
ignores maintenance hierarchy may return schedules that are numerically optimal
for the wrong problem: it may overlook a required reset, permit illegal
overlapping inspection windows, or certify a solution that fails under real
regulatory enforcement. By contrast, the A/B/C/D hierarchy captures the nested
reset logic assumed by this study and preserves the balance between operational
detail and tractability that is central to the thesis.
